\chapter{Concept Generation}
\label{ch:concept_generation}
\section{Session Overview}

Concept generation is the stage where your group explores the design space before deciding which ideas deserve deeper evaluation. The functions and solution options from Session \ref{sec:functional_decomposition_solution} become the building blocks for new concepts in this session.

By the end of this session, you should be able to:
\begin{enumerate}
    \item organise solution options using a concept classification tree and a concept combination table;
    \item combine sub-solutions into rough product concepts;
    \item describe each concept clearly enough for comparison; and
    \item shortlist several promising concepts for Session 11.
\end{enumerate}

\begin{tcolorbox}[title=Goal]
Today you will organise solution options into a classification tree and combination table, generate and describe at least five rough concepts by combining sub-solutions, and prepare to shortlist concepts for Session 11.
\end{tcolorbox}

\section{Key Tools}

\begin{itemize}
    \item \textbf{Concept classification tree}: a hierarchical way to group solution options for a sub-problem.
    \item \textbf{Concept combination table}: a matrix that helps the team combine one solution from each major sub-function.
    \item \textbf{Rough product concept}: an early concept description showing how several sub-solutions work together.
\end{itemize}

\section{Why These Tools Matter}

The purpose of concept generation is not to guess one perfect answer immediately. Instead, your team should use the functional decomposition from Session \ref{sec:functional_decomposition_solution} to generate several plausible combinations. This is particularly important in hydroponic products because different combinations of flow control, sensing, structure, and maintenance strategy can lead to very different concepts.

\section{Pre-Class Activity 1: Build a Classification Tree}

Start from the critical sub-functions identified in Session \ref{sec:functional_decomposition_solution} and group the available solutions into meaningful categories.

\subsection{Hydroponic Example}

For a hydroponic nutrient-delivery concept, a classification tree could be organised like this:
\begin{itemize}
    \item \textbf{Move solution}
    \begin{itemize}
        \item passive gravity feed
        \item electric pump
        \item peristaltic pump
    \end{itemize}
    \item \textbf{Monitor level}
    \begin{itemize}
        \item float indicator
        \item ultrasonic sensor
        \item conductivity or level probe
    \end{itemize}
    \item \textbf{Support structure}
    \begin{itemize}
        \item PVC frame
        \item aluminium frame
        \item modular plastic housing
    \end{itemize}
\end{itemize}

You may draw the tree as a diagram, or organise it as a structured list if that is easier.

\section{Pre-Class Activity 2: Create the Concept Combination Table}
\label{sec:create_concept_comb_table}
Use the classification work to create a table that combines solution options across the critical sub-functions.

\begin{table}[htbp]
\centering
\caption{Example concept combination table for a hydroponic monitoring product}
\label{tab:session10-combination-table}
\renewcommand{\arraystretch}{1.2}
\begin{tabularx}{\textwidth}{|>{\raggedright\arraybackslash}p{0.22\textwidth}|>{\raggedright\arraybackslash}p{0.22\textwidth}|>{\raggedright\arraybackslash}p{0.22\textwidth}|>{\raggedright\arraybackslash}X|}
\hline
\textbf{Move solution} & \textbf{Monitor level} & \textbf{Reduce leakage} & \textbf{Support structure} \\
\hline
Gravity feed & Float indicator & Sealed quick-connect fittings & PVC frame \\
\hline
Electric pump & Ultrasonic sensor & Gasketed joints & Aluminium frame \\
\hline
Peristaltic pump & Conductivity or level probe & Secondary containment tray & Modular plastic housing \\
\hline
\end{tabularx}
\end{table}

Each row is not automatically a finished concept. Instead, the table gives you a structured pool of options that can be recombined creatively.

\begin{tcolorbox}[title=Hint]
Group solutions by category and use the classification tree to ensure all options are considered before combining. The combination table is a menu, not a final answer, so mix and match creatively.
\end{tcolorbox}

\section{Pre-Class Activity 3: Generate Rough Concepts}

The objective of this activity is to generate rough product concepts by combining the solution concepts from different sub-problems within the concept combination table

\textbf{Procedure:}
\begin{enumerate}
    \item Refer to the output from activity in Section \ref{sec:create_concept_comb_table}, and generate rough product concepts by combining the solution concepts from different sub-problems within the concept combination table
    \begin{enumerate}
        \item Be creative and explore various combinations to generate a range of potential product concepts.
        \item Provide brief description to effectively convey the key aspects of the proposed rough product concepts.
        \item Provide information such as the technology utilized, working principles, intended functionality, and unique features of the concept.
    \end{enumerate}
    \item Repeat steps 1 for at least five generated concepts that show promise or align well with the project goals for further development.
\end{enumerate}

\begin{tcolorbox}[title=Checkpoint]
Generate at least five rough concepts. For each concept, record a name, selected sub-solutions, a short description of how it works, and at least one strength and weakness.
\end{tcolorbox}

% Use the combination table to generate at least five rough concepts.

% \subsection{Concept Description Template}

% For each concept, record:
% \begin{itemize}
%     \item a concept name;
%     \item the chosen sub-solutions;
%     \item a short description of how the concept works;
%     \item the main technologies involved; and
%     \item one or two likely strengths and weaknesses.
% \end{itemize}

% \begin{table}[htbp]
% \centering
% \caption{Template for documenting rough product concepts}
% \label{tab:session10-concept-template}
% \renewcommand{\arraystretch}{1.2}
% \begin{tabularx}{\textwidth}{|>{\raggedright\arraybackslash}p{0.14\textwidth}|>{\raggedright\arraybackslash}p{0.2\textwidth}|X|>{\raggedright\arraybackslash}p{0.2\textwidth}|}
% \hline
% \textbf{Concept name} & \textbf{Selected sub-solutions} & \textbf{How the concept works} & \textbf{Strengths and weaknesses} \\
% \hline
%  &  &  &  \\
% \hline
%  &  &  &  \\
% \hline
% \end{tabularx}
% \end{table}

\begin{tcolorbox}[title=Common Mistakes]
Do not prematurely pick a favourite concept or combine incompatible sub-solutions. Check whether combinations make sense in physics, maintenance, cost, and user needs before narrowing options.
\end{tcolorbox}

\section{In-Class Activity 1: Consolidate the Combination Table}

In class, compare the classification trees and combination tables prepared by each member. Combine the strongest ideas into one shared team version. Keep the discussion evidence-based: if one option appears more feasible, more innovative, or more strongly aligned with customer needs, explain why.

\section{In-Class Activity 2: Shortlist the Best Concepts}

After reviewing the full set of rough concepts, select the top five concepts for Session 11.

Use the following criteria for an initial shortlist:
\begin{enumerate}
    \item feasibility;
    \item alignment with customer needs and specifications;
    \item degree of innovation;
    \item integration of relevant physics principles; and
    \item potential for further development within the course.
\end{enumerate}

You do not need a full weighted scoring matrix yet, but you should record why each shortlisted concept survived the discussion.

\begin{tcolorbox}[title=Checkpoint]
When selecting concepts to carry forward, evaluate feasibility, alignment with customer needs and specifications, innovation level, integration of relevant physics principles, and development potential within the course. Document why each concept was shortlisted or discarded.
\end{tcolorbox}

\section{Summary and Next Steps}

At the end of this session, your group should have a shared classification tree, a shared combination table, and at least five clearly described rough concepts. These concepts will be screened and compared in Session \ref{ch:concept_selection}.
